% Constructive MTT Quantum Gravity III, corrected successor (v2)
\documentclass[11pt]{article}
\usepackage{series}

\hypersetup{
  unicode=true,
  pdftitle={Constructive MTT Quantum Gravity III: A Massive Scattering Benchmark and the Massless-Graviton Infrared Contract},
  pdfauthor={Peter Nero},
  pdfsubject={Corrected scattering and infrared boundary for the MTT quantum-gravity program}
}

\newcommand{\Has}{\mathcal H_{\mathrm{as}}}
\newcommand{\Hphys}{\mathcal H_{\mathrm{phys}}}
\newcommand{\TT}{\mathrm{TT}}
\newcommand{\OmIn}{\Omega_{-}}
\newcommand{\OmOut}{\Omega_{+}}

\title{Constructive MTT Quantum Gravity III:\\
\large A Massive Scattering Benchmark and the Massless-Graviton Infrared Contract}
\author{Peter Nero}
\date{July 2026\\Version 2}

\begin{document}
\maketitle

\begin{abstract}
This paper separates a massive or infrared-regulated scattering benchmark
from the physical massless-graviton problem in Modal Triplet Theory quantum
gravity. A positive transverse-traceless mass gap does not describe the
physical massless graviton, and finite-volume Gaussian and BRST results do
not by themselves construct the local Lorentzian theory required by
Haag--Ruelle scattering. The paper proves the operator statement that
for isometric incoming and outgoing wave operators, the scattering operator
is only a contraction in general and is unitary exactly when their ranges
coincide. It also proves that a normalized SPT factor leaves the leading
massless soft pole unchanged. Standard massive Haag--Ruelle theory is retained
as a conditional benchmark after its full axioms are supplied. The physical
gravity route instead requires a selected Lorentzian BRST theory, asymptotic
null geometry, soft-charge control, dressed states or inclusive observables,
infrared-regulator removal, and an appropriate completeness theorem. Current
MTT data do not yet provide that packet. The paper therefore supplies a
precise infrared dependency theorem and executable exit contract, not a
unitary MTT gravity S-matrix.
\end{abstract}

\section*{Version 2 Revision Note}
\begin{description}
\item[Supersedes:] Constructive MTT Quantum Gravity III, first release.
\item[Reason:] A positive TT gap, isometric wave operators, and SPT damping
were promoted beyond what they imply for massless gravitational scattering.
\item[Resolution:] Massive Haag--Ruelle theory is retained only as a
conditional benchmark; the exact wave-operator range criterion and SPT
infrared-neutrality result replace the former closure claim.
\item[Retained result:] Isometric wave operators give a contraction that is
unitary exactly when their ranges coincide.
\item[Remaining boundary:] Lorentzian BRST dynamics, null asymptotics, soft
charges, dressed or inclusive states, regulator removal, and completeness
remain open.
\end{description}

\tableofcontents

\section{Correction, purpose, and result map}

The first edition relied on the chain
\[
\begin{aligned}
 \text{SPT damping}
 &\Longrightarrow \text{gapped TT scattering}\\
 &\Longrightarrow \text{isometric wave operators}\\
 &\Longrightarrow \text{unitary gravitational }S\text{-matrix}.
\end{aligned}
\]
Every arrow was too strong. SPT is an ultraviolet spectral filter and does not
remove the massless soft pole. A positive TT gap changes the infrared theory.
The earlier papers did not construct the complete Lorentzian local net on
which scattering theory acts. Finally, isometric wave operators need not have
the same range, so their overlap need not be unitary.

This paper preserves the useful question but separates two branches:
\[
\boxed{\text{massive benchmark}}
\qquad\hbox{and}\qquad
\boxed{\text{physical massless-graviton contract}}.
\]
Section 3 identifies the errors in the old promotion. Section 4 states the
standard massive benchmark without claiming that MTT has supplied its
hypotheses. Section 5 proves the exact range criterion for scattering
unitarity. Section 6 replaces the massive shortcut by the correct soft-gravity
research contract. The final sections locate that contract in the present q79
program and record the remaining objects.

\subsection{The central picture in plain language}

Scattering compares what a system looks like in the far past with what it
looks like in the far future. The maps $\OmIn$ and $\OmOut$ embed those two
asymptotic descriptions into the interacting physical state space. Preserving
norms says that neither embedding loses probability. It does not say that the
two embeddings cover the same interacting states. Unitarity of the scattering
operator requires precisely that additional range statement.

For a massive local field, separated wave packets eventually move apart and
the Haag--Ruelle construction can exploit an isolated mass shell. Gravity has
a different infrared geometry. Its carrier is massless, its interaction is
long-ranged, and arbitrarily low-energy gravitons correlate the hard
particles with a soft cloud. The physical question is therefore not whether
an ultraviolet Gaussian factor makes a massive proof converge. It is whether
MTT selects the correct dressed or inclusive asymptotic object and proves its
infrared limits.

\section{What the current MTT results actually supply}

The current MTT quantum-gravity ledger contains several relevant but logically
earlier results:
\begin{enumerate}
\item a finite q79 internal TT operator and a positive free two-helicity block
at their declared finite-carrier tier;
\item a conditional reduction to the two-derivative Einstein/TEGR tensor
shape on an explicitly selected branch;
\item fixed-order quantum-GR effective-field-theory parity after the
renormalized action, Wilson coefficients, state, scale, gauge fixing, and
scheme are declared; and
\item corrected finite-volume SPT Gaussian and conditional BRST/BV interface
theorems.
\end{enumerate}

These are useful compatibility data. They do not yet emit a normalized
Lorentzian graviton field, a nonperturbative physical state space, a local
observable net, or asymptotic wave operators. In particular, the eigenvalues
of the finite internal TT Hessian are not automatically the invariant masses
of a Poincare representation. An explicit bridge from the internal carrier to
the four-dimensional Lorentzian spectrum is required before the word
\emph{mass gap} has its scattering-theory meaning.

The corrected predecessor papers also change what Part III may inherit.
Constructive QG I proves a bounded-volume Gaussian measure and a finite-mode
quartic Borel theorem, not an interacting gravity continuum. Constructive QG
II proves conditional BRST/BV interfaces, not a completed physical Hilbert
space. This paper therefore cannot begin by assuming those missing endpoints
have already been constructed.

\section{Why the old scattering proof fails}

\subsection{A finite slab is not an asymptotic region}

\begin{proposition}[Finite-slab obstruction]
Let a dynamics be defined only for $t\in[T_-,T_+]$ with both endpoints finite.
Then limits whose definition requires $t\to\pm\infty$ are not defined by that
finite-slab dynamics. A scattering construction needs either a global
dynamics or a controlled exhaustion in which the slab endpoints tend to
infinity.
\end{proposition}

\begin{proof}
The nets $t\mapsto U(t)$ and the comparison maps used in a wave-operator limit
must be evaluated for arbitrarily large positive or negative $t$. A function
defined only on a bounded interval has no such directed tail. Introducing a
sequence of larger slabs creates a second limit, whose existence and
independence must be proved rather than inferred from any one slab.
\end{proof}

The old definition of asymptotic flatness also included the existence of
Moller operators among the assumed consequences. Using that definition to
prove the same operators exist was circular. Geometric falloff, spectral
assumptions, propagation estimates, domains, and the wave-operator limits must
be stated as separate inputs.

\subsection{A positive TT gap is not the physical graviton}

The massless graviton has a lightlike one-particle shell. Assuming that the
physical spectrum starts above a strictly positive mass $\mu$ removes that
shell. Such an assumption can define a useful massive model, finite-volume
approximation, or infrared regulator, but its regulator-removal limit is an
additional theorem. It cannot be called the infrared completion of massless
gravity.

This distinction also prevents an ambiguity in Constructive QG I. Its
parameter $m^2>0$ makes the Euclidean bounded-volume covariance invertible. It
was explicitly introduced as an infrared shift, not selected as a physical
graviton mass.

\subsection{SPT smoothing does not remove the soft pole}

\begin{proposition}[Infrared neutrality of a normalized SPT filter]
\label{prop:ir-neutral}
Let $F:[0,\epsilon)\to\mathbb R$ satisfy
$F(u)=1+O(u)$ as $u\downarrow0$. Then
\[
 \frac{F(p^2)}{p^2}=\frac1{p^2}+O(1)
 \qquad (p^2\downarrow0).
\]
In particular, multiplying a massless propagator by such an SPT factor does
not remove its leading soft pole.
\end{proposition}

\begin{proof}
Write $F(u)=1+u r(u)$ with $r$ bounded near zero. Division by $u=p^2$
gives $F(u)/u=1/u+r(u)$.
\end{proof}

For the proper-time form
\[
 F(u)=\int_0^\infty e^{-tu}\,d\mu(t),
\]
normalization is $F(0)=\mu([0,\infty))=1$. If the first moment of $\mu$ is
finite, then $F(u)=1-u\int t\,d\mu(t)+o(u)$, so
Proposition~\ref{prop:ir-neutral} applies directly. The filter can strongly
change high momentum while leaving the leading infrared singularity intact.

\subsection{Euclidean decay is not Lorentzian locality}

Rapid decay of a Euclidean covariance does not by itself construct
microcausal Lorentzian observables. A nonpolynomial function of a Laplacian is
generally nonlocal. The old proof used the Gaussian multiplier both to improve
decay and to assume the strict locality needed by Haag--Ruelle theory. A
future construction must instead prove that its physical observable algebra
has the required commutation or almost-locality properties after continuation
and BRST reduction.

\section{The valid massive benchmark}

The corrected massive statement belongs to standard local quantum field
theory. It is retained because it gives a clean test that a future MTT
Lorentzian construction may attempt to satisfy in a genuinely massive sector.

\begin{theorem}[Conditional massive Haag--Ruelle benchmark]
\label{thm:massive}
Suppose a theory on Minkowski spacetime is supplied with a local,
translation-covariant observable net on a positive Hilbert space, a unique
vacuum, the spectrum condition, and an isolated stable one-particle mass
hyperboloid separated from the remaining spectrum as required by
Haag--Ruelle theory. Then the Haag--Ruelle limits define incoming and outgoing
multi-particle isometries from the corresponding asymptotic Fock space into
the physical Hilbert space. A unitary scattering operator follows only after
the relevant incoming and outgoing ranges are proved equal.
\end{theorem}

\begin{proof}
The existence and isometry statement is the standard Haag--Ruelle theorem
\cite{Ruelle1962,Haag1996}. The final sentence follows from
Theorem~\ref{thm:range} below.
\end{proof}

This is a conditional benchmark, not a new construction of the stated net.
It also does not apply directly to the massless graviton. Refined
Haag--Ruelle methods can construct stable massive particles in theories that
also contain massless excitations under additional regularity assumptions
\cite{Dybalski2005}; this still does not turn the graviton itself into an
isolated massive particle.

\section{What wave operators actually prove}

\begin{theorem}[Wave-operator range criterion]
\label{thm:range}
Let $\OmIn,\OmOut:\Has\to\Hphys$ be isometries and define
\[
 S=\OmOut^{*}\OmIn:\Has\to\Has.
\]
Then $\|S\|\leq1$. Moreover,
\[
 S^{*}S=I
 \quad\Longleftrightarrow\quad
 \Ran\OmIn\subseteq\Ran\OmOut,
\]
and
\[
 SS^{*}=I
 \quad\Longleftrightarrow\quad
 \Ran\OmOut\subseteq\Ran\OmIn.
\]
Consequently $S$ is unitary exactly when
$\Ran\OmIn=\Ran\OmOut$.
\end{theorem}

\begin{proof}
Let $P_\pm=\Omega_\pm\Omega_\pm^*$ be the orthogonal projections onto
the closed ranges of the two isometries. Then
\[
 S^*S=\OmIn^*P_+\OmIn,\qquad
 SS^*=\OmOut^*P_-\OmOut.
\]
Since orthogonal projections are contractions, so is $S$. The first
expression equals the identity precisely when $P_+\OmIn=\OmIn$, which is the
first range inclusion. The second is analogous. Both identities hold exactly
when both inclusions hold.
\end{proof}

\paragraph{Concrete foothold.}
Take $\Has=\mathbb C$, $\Hphys=\mathbb C^2$,
\[
 \OmIn z=(z,0),\qquad
 \OmOut z=(\cos\theta\,z,\sin\theta\,z).
\]
Both maps are isometries, but $S z=\cos\theta\,z$. Unless the two ranges
coincide, $S$ is not unitary. This one-dimensional example isolates the gap
in the first edition: norm preservation of each asymptotic embedding is not
asymptotic completeness.

For conventional scattering theory, wave operators are often described as
partial isometries on their initial subspaces. That terminology does not
repair the missing range theorem. In the present same-domain formulation,
$\OmOut^*\OmIn$ is guaranteed to be a contraction; it need not itself be a
partial isometry for arbitrary relative ranges.

\section{The physical massless-graviton route}

\subsection{Why ordinary Fock scattering is insufficient}

Long-range interactions do not switch off in the same way as short-range
massive interactions. In perturbative gravity, virtual and real soft
gravitons produce infrared structures already analyzed by Weinberg
\cite{Weinberg1965}. A naive hard-particle Fock state omits the correlated
soft gravitational field. Perturbative constructions therefore use either
inclusive observables, in which unresolved soft radiation is summed, or
appropriately dressed asymptotic states. Faddeev--Kulish methods originated
in QED \cite{KulishFaddeev1970}; a gravitational implementation constructs an
asymptotic dynamics and coherent soft clouds \cite{WareSaotomeAkhoury2013}.

These approaches are meaningful precedents, not results automatically owned
by MTT. They also carry choices and hypotheses concerning gauge invariance,
soft charges, asymptotic symmetries, factorization, and the class of
observables. Modern dressed-state analyses make explicit that hard and soft
sectors are correlated \cite{ChoiAkhoury2019}.

\subsection{Two legitimate output types}

A corrected MTT infrared program may target either of the following:
\begin{description}
\item[Inclusive output.] Infrared-safe transition probabilities or detector
observables after summing over experimentally unresolved soft gravitons, with
the resolution prescription and regulator removal declared.
\item[Dressed output.] A physical asymptotic state space whose hard states are
accompanied by the selected soft gravitational dressing, together with
well-defined incoming and outgoing maps and their range theorem.
\end{description}
The two descriptions may agree on suitable observables, but they are not
identical data structures. A paper must declare which one it constructs.

\subsection{Selected MTT infrared exit contract}

For a physical MTT gravity scattering theorem, the following objects must come
from one compatible source chain:
\begin{enumerate}
\item \textbf{Lorentzian theory.} A selected gauge-invariant action or
observable net, a gravitational BRST/BV construction, a physical state, and
controlled continuation from any Euclidean input.
\item \textbf{Massless spectrum.} A normalized two-helicity massless graviton
sector with a proved map from the finite q79 TT carrier to the Lorentzian
Poincare or asymptotic representation.
\item \textbf{Asymptotic geometry.} A declared asymptotically flat spacetime
or null-infinity structure, with dynamics defined for arbitrarily large
times. Falloff assumptions must not include the desired wave operators.
\item \textbf{Soft data.} Soft charges, memory sectors, and the action of
large gauge or asymptotic symmetries on the physical observables.
\item \textbf{Infrared prescription.} Either an inclusive resolution map or a
selected dressed-state construction, including gauge/BRST compatibility.
\item \textbf{Removal theorem.} Independence of auxiliary graviton masses,
finite volumes, adiabatic switches, and soft-energy cutoffs in the declared
observable topology.
\item \textbf{Scattering ranges.} Incoming and outgoing maps, the range
relations of Theorem~\ref{thm:range}, and the chosen completeness statement.
\end{enumerate}

The first edition did not supply these seven objects. Importantly, SPT damping
does not close item 4 or 5: by Proposition~\ref{prop:ir-neutral}, its normalized
ultraviolet factor leaves the leading soft pole unchanged.

\section{Relation to the q79 heterotic route}

The selected compatible ultraviolet candidate in the current MTT program is
the q79 Fu--Yau heterotic branch, not the SPT filter by itself. That route
could eventually constrain the massless spectrum, BRST data, and interactions
entering item 1 and item 2 of the exit contract. At present its worldsheet
contract has five rows available, two partial, and five open. In particular,
the exact worldsheet theory, complete BV vertices, tadpole/vacuum control,
infrared and soft completion, and all-genus/nonperturbative definition are not
yet available.

Fixed-genus string amplitudes also do not automatically solve the
four-dimensional asymptotic problem. One must identify the physical
four-dimensional states, take the relevant compactification and low-energy
limits, and establish the soft/inclusive or dressed prescription in that
sector. The heterotic route is therefore compatible with the contract above,
but it has not yet discharged it.

\section{Status ledger}

\begin{center}
\small
\begin{tabular}{p{0.28\textwidth}p{0.20\textwidth}p{0.43\textwidth}}
\hline
Object & Status & Meaning \\
\hline
Massive Haag--Ruelle theorem & Imported benchmark & Applies after a local positive theory with the required isolated mass shell is supplied. \\
Finite q79 TT carrier & Available at finite tier & Not yet a Lorentzian mass-shell or asymptotic-state theorem. \\
SPT finite-volume Gaussian & Available conditionally & Controls a Euclidean ultraviolet model; leaves the leading massless soft pole intact. \\
Gravitational BRST/BV state space & Open & Corrected QG II gives interfaces, not the selected quantum measure or physical completion. \\
Massless graviton dressing or inclusive map & Open & No selected MTT soft-sector operator has been emitted. \\
IR-regulator removal & Open & No uniform physical limit for mass, volume, switching, and soft cutoffs is proved. \\
Wave-operator range equality & Open & Isometries alone do not imply a unitary scattering operator. \\
Physical MTT gravity $S$-matrix & Open & Requires the complete seven-object contract. \\
\hline
\end{tabular}
\normalsize
\end{center}

\section{Version delta}

Relative to version 1, this successor:
\begin{itemize}
\item withdraws the claim that the three constructive papers complete
nonperturbative MTT quantum gravity;
\item reclassifies the positive-gap argument as a massive or
infrared-regulated benchmark;
\item removes the circular definition that assumed Moller operators inside
asymptotic flatness;
\item distinguishes finite slabs from the infinite-time limit needed for
scattering;
\item proves the exact wave-operator range criterion and corrects the claim
that isometries alone make the $S$-matrix unitary;
\item proves that a normalized SPT factor leaves the massless soft pole
unchanged;
\item replaces the naive graviton Fock-space conclusion by inclusive and
dressed alternatives; and
\item states the seven-object MTT infrared exit contract.
\end{itemize}

\section{Conclusion}

The corrected result is sharp. Standard massive Haag--Ruelle theory remains a
valuable benchmark, but a positive graviton mass gap is not the infrared
physics of our universe. SPT smoothing is ultraviolet data and, when
normalized at zero momentum, does not cure the massless soft singularity.
Even after incoming and outgoing isometries exist, unitarity requires equality
of their ranges.

What MTT has achieved here is therefore not a completed gravity $S$-matrix.
It has a finite TT carrier, a conditional Einstein reduction, fixed-order EFT
parity, and now a corrected map of the infrared problem. The next genuine
advance is to construct one selected massless soft-sector packet: Lorentzian
physical states, asymptotic charges, a dressing or inclusive prescription, and
its regulator-removal and range certificates. That packet would connect the
current finite q79 structure to physical gravitational scattering without
hiding the long-range problem behind an artificial mass gap.

% BEGIN MTT MANAGED COMPUTATIONAL EVIDENCE
\section*{Computational Evidence and Reproducibility}
The massive benchmark and massless infrared obligations are separated within this paper. The mapped open strict-upgrade row does not close the massless-graviton limit and is cited only as an adjacent unresolved source obligation.

The referenced rows are frozen to the curated results repository state identified below. Hashes are grouped in eight-character blocks for line breaking.
\begin{quote}\small
\textbf{Repository:} \url{https://github.com/PeterNero/mtt-results-repro}\\
\textbf{Commit:} \texttt{31247ebb 5c22f3fb b5443024 365433c6 ee0bff4a}\\
\textbf{Manifest:} \path{release/result_manifest.json}\\
\textbf{Manifest SHA-256:}\\
\texttt{fb399689 60b00584 631dbf53 1a708e18}\\
\texttt{ef928d6b 6d935119 c185d7f6 32b1e7cd}
\end{quote}

Tier labels are quoted verbatim from that manifest. A row used directly supports only the specific computational statement identified above; a corpus-state cross-check does not prove this paper's local theorems; and an open row is evidence of an unresolved obligation, never of closure.

\par\begingroup\dimen0=\pagegoal\advance\dimen0 by -\pagetotal\ifdim\dimen0<7\baselineskip\newpage\fi\endgroup
\paragraph{Open boundary (not evidence of closure).}
\begin{itemize}
\setlength{\itemsep}{0.25em}
\item \path{A05/strict_upgrade_ledger} (\emph{open}).\par Current 2/9 strict no-knob upgrade ledger.
\end{itemize}

No imported row changes theorem ownership or promotes a neighboring claim: all local statements retain their stated hypotheses, domains, and limitations.
% END MTT MANAGED COMPUTATIONAL EVIDENCE

\begin{thebibliography}{99}

\bibitem{Ruelle1962}
D.~Ruelle,
\emph{On the asymptotic condition in quantum field theory},
Helv.\ Phys.\ Acta \textbf{35} (1962), 147--163.

\bibitem{Haag1996}
R.~Haag,
\emph{Local Quantum Physics},
2nd ed., Springer, 1996.

\bibitem{Dybalski2005}
W.~Dybalski,
\emph{Haag--Ruelle scattering theory in presence of massless particles},
Lett.\ Math.\ Phys.\ \textbf{72} (2005), 27--38,
doi:10.1007/s11005-005-2294-6.

\bibitem{Weinberg1965}
S.~Weinberg,
\emph{Infrared photons and gravitons},
Phys.\ Rev.\ \textbf{140} (1965), B516--B524,
doi:10.1103/PhysRev.140.B516.

\bibitem{KulishFaddeev1970}
P.~P.~Kulish and L.~D.~Faddeev,
\emph{Asymptotic conditions and infrared divergences in quantum
electrodynamics},
Theor.\ Math.\ Phys.\ \textbf{4} (1970), 745--757,
doi:10.1007/BF01066485.

\bibitem{WareSaotomeAkhoury2013}
J.~Ware, R.~Saotome, and R.~Akhoury,
\emph{Construction of an asymptotic S matrix for perturbative quantum
gravity},
JHEP \textbf{10} (2013), 159,
doi:10.1007/JHEP10(2013)159,
arXiv:1308.6285.

\bibitem{ChoiAkhoury2019}
S.~Choi and R.~Akhoury,
\emph{Subleading soft dressings of asymptotic states in QED and
perturbative quantum gravity},
JHEP \textbf{09} (2019), 031,
arXiv:1907.05438.

\end{thebibliography}

\end{document}
